\chapter*{Lampiran}
\addcontentsline{toc}{chapter}{Lampiran}

% ============================================================
% Lampiran A: Rubrik Penilaian
% ============================================================
\section*{Lampiran A: Rubrik Penilaian Tugas Praktik Pemrograman Game}

\begin{table}[h]
\centering
\small
\begin{tabular}{|>{\raggedright\arraybackslash}p{3cm}|>{\raggedright\arraybackslash}p{3.5cm}|>{\raggedright\arraybackslash}p{3.5cm}|>{\raggedright\arraybackslash}p{3.5cm}|}
\hline
\textbf{Kriteria} & \textbf{Sangat Baik} & \textbf{Baik} & \textbf{Perlu Perbaikan} \\
\hline
Implementasi Unity & Semua fitur berfungsi, struktur rapi & Sebagian fitur berfungsi & Banyak fitur tidak berfungsi \\
\hline
Kode C\# & Bersih, terorganisir, komentar & Cukup rapi & Tidak terorganisir \\
\hline
Gameplay & Mechanic lengkap, balance & Mechanic dasar ada & Mechanic tidak lengkap \\
\hline
UI/UX & Interface jelas, feedback & Cukup jelas & Kurang informatif \\
\hline
Audio & SFX dan BGM sesuai & Ada audio dasar & Tidak ada/tidak sesuai \\
\hline
Build & Berhasil build, jalan stabil & Build berhasil & Build gagal \\
\hline
\end{tabular}
\caption{Rubrik Penilaian Tugas Praktik Pemrograman Game}
\end{table}

% ============================================================
% Lampiran B: Contoh Template Laporan Proyek Game
% ============================================================
\section*{Lampiran B: Contoh Template Laporan Proyek Game}
\begin{enumerate}
  \item Judul dan Identitas Mahasiswa
  \item Deskripsi Game (konsep, GDD singkat)
  \item Screenshot dan fitur gameplay
  \item Implementasi (cuplikan kode penting)
  \item Hasil Build dan Testing
  \item Refleksi dan Kesimpulan
\end{enumerate}

% ============================================================
% Lampiran C: Glosarium Istilah Game Development
% ============================================================
\section*{Lampiran C: Glosarium Istilah Game Development}
\begin{itemize}
  \item \textbf{GameObject}: Objek dasar dalam Unity yang dapat ditempatkan di scene
  \item \textbf{Component}: Bagian fungsional yang ditambahkan ke GameObject
  \item \textbf{Prefab}: Template GameObject yang dapat diinstansiasi berulang
  \item \textbf{Scene}: Container untuk level atau bagian game
  \item \textbf{MonoBehaviour}: Kelas dasar untuk skrip Unity
  \item \textbf{Rigidbody}: Komponen fisika untuk gerakan
  \item \textbf{Collider}: Komponen untuk deteksi tabrakan
  \item \textbf{Animator}: Komponen untuk mengontrol animasi
  \item \textbf{Canvas}: Area untuk elemen UI
  \item \textbf{Instantiate}: Membuat instance baru dari Prefab
  \item \textbf{MDA}: Mechanics, Dynamics, Aesthetics framework
  \item \textbf{GDD}: Game Design Document
\end{itemize}

% ============================================================
% Lampiran D: Referensi Tambahan
% ============================================================
\section*{Lampiran D: Referensi Tambahan}
\begin{itemize}
  \item Unity Manual: \url{https://docs.unity3d.com/Manual/}
  \item Unity Scripting API: \url{https://docs.unity3d.com/ScriptReference/}
  \item Unity Learn: \url{https://learn.unity.com/}
  \item C\# Documentation: \url{https://learn.microsoft.com/dotnet/csharp/}
  \item Game Programming Patterns: \url{https://gameprogrammingpatterns.com/}
\end{itemize}
